\scanpage{105}
\mspage{54}
J'ai cerclé les notions qui ont été vraiment utiles. Pratiquement, on peut parier que c'est la notion la plus forte, $W(4a)$, qui est la seule importante.
Voici les implications tautologiques (à l'exception de $W(4a')\Longrightarrow W(4b)$, qui est \deleted{\unsure{la principale}} \inserted{\unsure{démontré|démontrée}} dans la dém. de la prop. 4 (pages 6–10).
% slanted
\margin{Cette démonstration utilise une notion plus forte de $W$-fibration que celle du th. 3 p. 51}
\diagram{\begin{tikzcd}[ampersand replacement=\&, row sep=small,column sep=small]
\& a\arrow[d,Rightarrow] \& {} \\
\& a'\arrow[d,Rightarrow] \& {} \\
\& b\arrow[dl,Rightarrow]\arrow[d,Rightarrow]\arrow[dr,Rightarrow] \& {} \\
b' \& c\arrow[d,Rightarrow] \& b'' \\
\& c' \& {}
\end{tikzcd}}
% slanted
\margin{implication douteuse il faut retoucher un peu b), en renforçant la notion de $W$-fibration utilisée pour énoncer b)}

J'ai surtout sué pour montrer que $c'\Rightarrow c$, ou $c\Longrightarrow b$, sans y arriver. L'effort m'a permis de me débarrasser de diverses idées fausses, et de les remplacer par des justes, cf plus bas. Il y a quand même un résultat positif qui surnage :

\uline{Proposition 1} : La condition $b$ (et sans doute même $b''{}^*$, mais je ne vais pas le prouver \unsure{\ldots}) implique \inserted{sans doute} $a$ pour $\bar W$, le saturé fort de $W$. \deleted{\unsure{D}} \inserted{En tous cas,} si $W$ est fortement saturé, les conditions $a,a',b$ (et sûrement aussi $b'$ et $b''$) sont équivalentes.
% slanted
\margin{et en supposant, dans le cas de $b''$, que $W$ satisfait la condition \uline{$b'$} des Limites inductives $W(6\,\mathrm{bis})$. cf. \deleted{le précédent} prop. 3 de ``\unsure{variations du|variations sur le|corollaire de} \unsure{\ldots}'' page 14}

Démonstration de $b\Longrightarrow a$, \deleted{\unsure{pour $W$}} \inserted{quand $W$ est ft saturé.} \deleted{\unsure{Il}} En fait, moyennant b) (et sans $=$ supposer $W$ ft \unsure{saturé|saturée}) nous allons prouver que si on a
\diagram{\begin{tikzcd}[ampersand replacement=\&]
X\arrow[d,"f"'] \& X'\arrow[l,"q"']\arrow[d,"f'"] \\
Y \& Y'\arrow[l,"p"]
\end{tikzcd}}
cartésien dans $\mathrm{Cat}$, avec
% archive sheet number: 105
